\section{Introduction}

Finite-amplitude perturbations probe a form of stability that is not captured by the response to infinitesimal deviations. Even when a state is locally stable, a sufficiently large disturbance may redirect its evolution toward a different macroscopic outcome. This distinction is central to ecological resilience \cite{holling1973resilience,scheffer2001catastrophic} and to basin-based measures of nonlinear stability \cite{menck2013basin}. In spatially extended systems, perturbation magnitude is not the only relevant quantity: the initial arrangement of disturbed sites also determines where interfaces form and how they can propagate through the surrounding medium.

Spatial organization connects local interactions to collective dynamics \cite{levin1992pattern,durrett1994spatial}. Pattern formation and front propagation provide useful descriptions of how interfaces emerge, move, and reorganize extended structures \cite{cross1993patterns,vansaarloos2003fronts}. In ecological systems, feedback between populations and transported resources can reshape spatial patterns and alter the response to perturbations \cite{rietkerk2004patchiness,bel2012gradual,rietkerk2021tipping}. Spatial evolutionary games offer a complementary microscopic setting in which the competing domains correspond to different strategies. Local imitation and interaction structure generate correlations and persistent spatial patterns \cite{nowak1992spatial,nowak1994spatial}, although spatial organization does not universally favor cooperation \cite{hauert2004spatial}. Results on graph-structured populations and network reciprocity further show how interaction topology affects selection \cite{nowak2006rules,ohtsuki2006rule,allen2017structure}, connecting spatial game dynamics to nonequilibrium statistical physics \cite{szabo2007graphs,roca2009review,perc2017cooperation,szabo2002volunteering}.

Environmental feedback adds a second evolving field to this picture \cite{perc2010coevolutionary}. Ecological public-goods models have shown that dispersal and population feedback can generate spatial structure \cite{wakano2009spatial}. Game--environment feedback can also drive depletion, recovery, and oscillatory dynamics, with outcomes depending on heterogeneity and relative time scales \cite{weitz2016oscillating,hauert2019asymmetric,tilman2020feedback}. In stochastic games and games with changing interaction conditions, strategies modify the environment in which subsequent competition takes place \cite{hilbe2018stochastic,su2019transitions}. When these effects are local, strategy and resource distributions can form coupled spatial patterns \cite{lin2019spatial,chen2025localized,zhong2026spatial}. The resource field is therefore not merely a background parameter: it records and responds to the spatial history of the competing populations.

This coupling raises a question that cannot be resolved from the statistics of either field alone. Can two states with the same strategy configuration and the same resource-value distribution and autocorrelation have different probabilities of collapse because the resource pattern is positioned differently relative to the strategy interface? Comparing independently evolved configurations does not isolate this effect, because interface geometry, resource abundance, and spatial correlations generally change together. Translating only the resource field on a periodic lattice provides a controlled comparison: its distribution and autocorrelation remain unchanged, as does the entire strategy configuration, while the relative alignment of the two fields changes.

Here we investigate finite disturbances in a stochastic lattice evolutionary game coupled to resource renewal, depletion, and diffusion. We first compare compact and dispersed perturbations with identical defector counts and initial resource loss to isolate the effect of disturbance geometry. We then translate the resource field of grown disturbances to test the role of relative alignment in subsequent collapse. Spatial measurements and first-passage probabilities connect the geometric contrast to disturbance establishment and the alignment contrast to interfacial resource structure. Exact first-visit, resource-balance, and sensitivity relations identify the microscopic terms underlying these observations. Finally, resource permutations along a fixed boundary distinguish immediate interfacial responses from their subsequent evolution. Throughout, collapse and recovery are defined by finite-time first-passage criteria \cite{redner2001firstpassage}, rather than by an inferred asymptotic phase boundary.
